% !TeX root = ../papers/nc-p3.tex
\section{被約四頂点切断の六つの反対座標積}
\label{sec:finite-marked-tetrahedral-products}
\begin{theorem}\label{thm:finite-marked-tetrahedral-products}
共通の印付き切断で $\Sigma$ が四つの被約点とし、印の接次元は1以下とする。
偶位置の座標を四頂点に対応する $a_i,b_i$ ($1\le i\le4$) に選び、
二乗空間を一致させる単元により自明化すると
\[
                      b_i b_j=a_i a_j \quad(i\ne j)
\]
となる。偶位置が対称でなければ六つの $g_{ij}=a_i b_j-a_j b_i$ ($i<j$) は独立で、
その各支持スキームは長さ1または2であり、異なる六つの支持は交わらない。
$d_i=b_i^2-a_i^2$ を $a$ の一意な二次式と見ると、どの三つの異なる添字でも
$d_i a_j a_k=0$ が点環で成り立つ。少なくとも三つの $d_i$ は非零である。
\end{theorem}
\begin{proof}
長さ4の積跡模型を使う。四つの階数一行列を $u_i u_i^t$ とし、
三次元の $u_i$ はどの三個も独立である。
$E$ の二つの評価を $A=\sum b_i u_i u_i^t$, $B=\sum a_i u_i u_i^t$ とする。
局所自明化による各列全体への単元倍は以下の消失条件に影響しない。
位置1の三射影座標を $l$、$l_i=u_i^tl$ と書く。位置3にも同じ座標表示を使う。
例外鉛筆は二頂点の張る直線で、Artin基底上でもそのような鉛筆は一定である。
実際相対的な次数2の階数一因子は定数な四点の直和因子二つに限られる。
鉛筆 $ij$ に対応する位置1の直線は $l_i=l_j=0$ である。

印以外の任意の局所因子で $A,B$ が比例しないとき、前節までの縮約行列の主元消去により、
正規化した差は $z V$ となる。$S/\operatorname{ann}(z)$ 上で $A,V$ は剰余で独立であり、
二点 $i,j$ の張る固定鉛筆の基底となる。従って外積 $A\wedge B$ の $ij$ 以外の成分は零、
$ij$ 成分は $a_k,b_k$ ($k\notin\{i,j\}$) および $l_i,l_j$ に消される。
印でも、$S=k[t]/(t^m)$, $z=t^s$ のとき $l/t$ を用いて
$t^{m-s-1}$ を法として同じ鉛筆を得る。従って全ての異なる $i,j,k$ で $b_k g_{ij}$ は
印に支持されるsocleに入り、$l_i g_{ij}=l_j g_{ij}=0$ が全域で成り立つ。
実際、指定鉛筆以外の外積成分が印で残ってもsocleに属し、$l_i$ は印で零である。
$m-s\le1$ なら初めから全外積成分がsocleに入る。被約の印でも結論は成立する。

単元 $h$ を $W^2=h U^2$ と取り、$P_{ij}(a)$ を $h^{-1} b_i b_j$ の二次持上げとする。
$x$ を印の一つのsocle生成元とし、$C(a)$ をそれを表す一意な三次式とする。
各相異なる $i,j,k$ について
\[
 a_i P_{jk}-a_j P_{ik}=c_{ijk} C
\]
が多項式恒等式である。$k$ を固定し、他の三添字 $i,j,n$ の三つの式に
$a_n,a_i,a_j$ を掛けて巡回和を取ると左辺は零である。
$C\ne0$ と三一次式の独立性から全係数が零である。
従って $P_{ik}=a_i h_k$ ($i\ne k$) と書け、$P$ の対称性から $h_k=c a_k$ が
共通の非零定数 $c$ で成立する。$ch$ の平方根で $b$ の自明化を変えれば最初の式を得る。
これで $a_k g_{ij}=b_k g_{ij}=0$ は印を含む全域で成り立つ。
偶側が非対称なら定理\ref{thm:finite-opposite-product-gap}より六元は独立で全て非零である。

各 $g_{ij}$ の支持は偶位置の座標直線 $ij$、位置1の直線 $l_i=l_j=0$ の積で
六関係を満たす。従って一般の二式の交点数から長さは2以下である。
印以外の局所因子上では既に高々一つの割線に限られる。
印では $A(p)=B(p)$ の階数は $3-\dim T_pX\ge2$ なので、これは高々一本の座標直線に入る
（異なる直線の交点は階数一頂点である）。偶直線条件が全域で成り立つから印でも一つに限られる。
最後に $d_i a_j=-b_i g_{ij}$ ($i\ne j$、外積は添字交換で符号を変える) より四次の消失を得る。
二つの $d_i,d_j$ がともに零なら $b_i,b_j$ は $g_{ij}$ をともに消す。
その支持は $b$ 側でも直線 $ij$ に入るので空となり、独立性に反する。
\end{proof}
